% Options for packages loaded elsewhere
\PassOptionsToPackage{unicode}{hyperref}
\PassOptionsToPackage{hyphens}{url}
\PassOptionsToPackage{dvipsnames,svgnames,x11names}{xcolor}
\documentclass[
  11pt,
  a4paper,
]{article}
\usepackage{xcolor}
\usepackage[margin=2.6cm]{geometry}
\usepackage{amsmath,amssymb}
\setcounter{secnumdepth}{5}
\usepackage{iftex}
\ifPDFTeX
  \usepackage[T1]{fontenc}
  \usepackage[utf8]{inputenc}
  \usepackage{textcomp} % provide euro and other symbols
\else % if luatex or xetex
  \usepackage{unicode-math} % this also loads fontspec
  \defaultfontfeatures{Scale=MatchLowercase}
  \defaultfontfeatures[\rmfamily]{Ligatures=TeX,Scale=1}
\fi
\usepackage{lmodern}
\ifPDFTeX\else
  % xetex/luatex font selection
  \setmainfont[]{Libertinus Serif}
  \setmonofont[Scale=0.85]{Latin Modern Mono}
  \setmathfont[]{Latin Modern Math}
\fi
% Use upquote if available, for straight quotes in verbatim environments
\IfFileExists{upquote.sty}{\usepackage{upquote}}{}
\IfFileExists{microtype.sty}{% use microtype if available
  \usepackage[]{microtype}
  \UseMicrotypeSet[protrusion]{basicmath} % disable protrusion for tt fonts
}{}
\makeatletter
\@ifundefined{KOMAClassName}{% if non-KOMA class
  \IfFileExists{parskip.sty}{%
    \usepackage{parskip}
  }{% else
    \setlength{\parindent}{0pt}
    \setlength{\parskip}{6pt plus 2pt minus 1pt}}
}{% if KOMA class
  \KOMAoptions{parskip=half}}
\makeatother
\usepackage{longtable,booktabs,array}
\newcounter{none} % for unnumbered tables
\usepackage{calc} % for calculating minipage widths
% Correct order of tables after \paragraph or \subparagraph
\usepackage{etoolbox}
\makeatletter
\patchcmd\longtable{\par}{\if@noskipsec\mbox{}\fi\par}{}{}
\makeatother
% Allow footnotes in longtable head/foot
\IfFileExists{footnotehyper.sty}{\usepackage{footnotehyper}}{\usepackage{footnote}}
\makesavenoteenv{longtable}
\setlength{\emergencystretch}{3em} % prevent overfull lines
\providecommand{\tightlist}{%
  \setlength{\itemsep}{0pt}\setlength{\parskip}{0pt}}

\providecommand{\E}{\mathbb{E}}
\providecommand{\N}{\mathbb{N}}
\providecommand{\Z}{\mathbb{Z}}
\providecommand{\R}{\mathbb{R}}
\providecommand{\eps}{\varepsilon}
\providecommand{\ceil}[1]{\left\lceil #1\right\rceil}
\providecommand{\floor}[1]{\left\lfloor #1\right\rfloor}
\AtBeginDocument{\let\setminus\smallsetminus}
\usepackage[most]{tcolorbox}
\newtcolorbox{warning}{colback=red!5,colframe=red!65!black,boxrule=0.8pt,arc=2pt,left=6pt,right=6pt,top=5pt,bottom=5pt}
\usepackage{etoolbox}
\AtBeginEnvironment{longtable}{\small}
\setlength{\emergencystretch}{3em}
\usepackage{fancyhdr}
\pagestyle{fancy}
\fancyhf{}
\fancyhead[L]{\small\itshape Candidate result 2211.01032\_\_03 --- GPT-6 Astra, awaiting review}
\fancyhead[R]{\small\thepage}
\renewcommand{\headrulewidth}{0.2pt}
\usepackage{bookmark}
\IfFileExists{xurl.sty}{\usepackage{xurl}}{} % add URL line breaks if available
\urlstyle{same}
\hypersetup{
  pdftitle={A candidate proof of Conjecture 9.5 in Random Embeddings of Graphs: The Expected Number of Faces in Most Graph\ldots{}},
  pdfauthor={github.com/graph-theory-AI --- machine-generated candidate writeup},
  colorlinks=true,
  linkcolor={blue!50!black},
  filecolor={Maroon},
  citecolor={Blue},
  urlcolor={blue!50!black},
  pdfcreator={LaTeX via pandoc}}

\title{A candidate proof of Conjecture 9.5 in Random Embeddings of
Graphs: The Expected Number of Faces in Most Graph\ldots{}}
\author{\href{https://github.com/graph-theory-AI}{github.com/graph-theory-AI}
--- machine-generated candidate writeup}
\date{16 September 2026}

\begin{document}
\maketitle
\begin{abstract}
A fixed-vertex decomposition proves E{[}F\_n{]} \textless= ln n + O(1),
with the stronger upper bound ln n + 1/4 + gamma/2 + o(1). This
candidate proof was generated by GPT-6 Astra in a single model pass for
the Graph Theory AI project. It was selected from the campaign because
the model marked it \texttt{would\_publish:\ true}; it has not been
checked by the adversarial referee, a human mathematician, or a novelty
search.
\end{abstract}

\section{The problem and the claimed
result}\label{the-problem-and-the-claimed-result}

This document concerns
\href{https://graph-theory-ai.github.io/graph-conjectures/arxiv/2211.01032__03/}{Conjecture
9.5} from \emph{Random Embeddings of Graphs: The Expected Number of
Faces in Most Graph\ldots{}}.

\subsection{Problem statement}\label{problem-statement}

Conjecture. The expected number of faces in a non-orientable random
embedding of the complete graph \(K_{n}\) is at most \(\ln(n) + O(1)\).

\subsection{Catalog context}\label{catalog-context}

The authors extend their investigation to non-orientable surfaces, where
a random embedding is defined by choosing a random rotation system
together with a random signature assigning each edge a sign
independently with probability \(1/2\). Computer simulations support a
logarithmic behaviour analogous to the orientable case.

\subsection{Claimed result}\label{claimed-result}

\textbf{Model verdict: proved} (confidence: high).

A fixed-vertex decomposition proves E{[}F\_n{]} \textless= ln n + O(1),
with the stronger upper bound ln n + 1/4 + gamma/2 + o(1).

\emph{The model's own caveat:} The proof uses exactly the stated
uniform-rotation, independent-fair-signature model; literature priority
has not been checked.

\section{Model writeup}\label{model-writeup}

The text below is the model output verbatim, apart from moving its
machine-readable verdict block into the summary above and shifting
Markdown heading levels for this document. Mechanical TeX defects and
equation tags were normalized where needed for compilation.

\subsection{A fixed-vertex proof}\label{a-fixed-vertex-proof}

Let \(F_n\) denote the number of faces in the random signed-rotation
embedding of \(K_n\) specified in the question.

The key is to separate faces meeting one fixed vertex from faces
avoiding it. Each class has expected size at most \(\frac12\ln n+O(1)\).
For the avoiding faces, an unvisited vertex supplies an escape option at
every exploration step. This controls local congestion and removes the
long-face obstruction in the previous attempt.

The flag-matching representation from that attempt is valid and is
justified below. Its short-face enumeration is not needed.

\subsubsection{Theorem}\label{theorem}

As \(n\to\infty\), \[
\boxed{
\mathbb E F_n
\le
\ln n+\frac14+\frac{\gamma}{2}
+O\!\left(\frac{(\ln n)^2}{n}\right),
}
\] where \(\gamma\) is Euler's constant. In particular, \[
\mathbb E F_n\le \ln n+O(1).
\]

All logarithms below are natural unless a base is indicated.

\begin{center}\rule{0.5\linewidth}{0.5pt}\end{center}

\subsubsection{1. The matching model and its exposure
rule}\label{the-matching-model-and-its-exposure-rule}

Put \(d=n-1\). For each incidence of an edge \(uv\) at \(u\), introduce
two flags \[
(u,v,0),\qquad (u,v,1).
\] Thus there are \(2d\) flags at each vertex and \(2n(n-1)\) flags
altogether.

Use the following matchings:

\begin{itemize}
\tightlist
\item
  \(J_u\) pairs the two flags of each incidence at \(u\).
\item
  \(A\) pairs corresponding flags at opposite ends of an edge: \[
  (u,v,s)\longleftrightarrow(v,u,s).
  \]
\item
  \(V_u\) is a matching on the flags at \(u\), describing the arcs
  between consecutive edge attachments around its vertex-disc.
\end{itemize}

The boundary components of the resulting ribbon graph are exactly the
connected components of \[
A\cup V,\qquad V=\bigcup_u V_u.
\] A face of length \(k\) contains \(k\) matching edges from \(A\),
\(k\) from \(V\), and \(2k\) flags. Consequently, \[
\sum_{\text{faces }f}|f|=n(n-1).
\]

\paragraph{Distribution of the local
matchings}\label{distribution-of-the-local-matchings}

Let \[
\mathcal H_d
=
\{W:\ W\text{ is a perfect matching and }J_u\cup W
       \text{ is one alternating }2d\text{-cycle}\}.
\] In the stated random-embedding model, we may take the \(V_u\) to be
independent and uniform on \(\mathcal H_d\), with \(A\) fixed as above.

Here is a verification. Choose independent fair bits \(b_{u,e}\) at all
incidences, and define the signature of \(e=uv\) by \[
\varepsilon_e=b_{u,e}\oplus b_{v,e},
\] with an immaterial fixed adjustment if the opposite sign convention
is used. These signatures are independent and fair. Relabelling the
flags at incidence \((u,e)\) using \(b_{u,e}\) makes the edge matching
canonical.

At a vertex, a cyclic order and the \(d\) incidence bits specify an
element of \(\mathcal H_d\). Every element has exactly two such
representations, corresponding to the two directions around its
alternating cycle. Hence \[
|\mathcal H_d|=2^{d-1}(d-1)!,
\] and the local matchings are independent and uniform, as asserted.

\paragraph{Local exposure lemma}\label{local-exposure-lemma}

Suppose \(r<d\) pairs of \(V_u\) have been exposed, forming a matching
\(Q_u\).

Because \(J_u\cup V_u\) is one cycle, \(J_u\cup Q_u\) consists of
\(d-r\) alternating paths. Its endpoints are exactly the flags whose
\(V_u\)-mates are still unexposed.

If \(r\le d-2\), then, for any such flag \(x\), its mate is uniform
among \[
2(d-r-1)
\qquad\text{(1)}
\] possible flags: all the other free flags except the opposite endpoint
of the same \(J_u\cup Q_u\)-path. If \(r=d-1\), the last pair is forced.

Indeed, contracting the alternating paths reduces the conditional
completion problem to a uniform element of \(\mathcal H_{d-r}\). Joining
two different paths is allowed, while closing a proper local cycle is
forbidden. All allowed partners have the same number of completions.

This rule remains valid under adaptive exposure using only previously
revealed information.

\begin{center}\rule{0.5\linewidth}{0.5pt}\end{center}

\subsubsection{2. Faces meeting a fixed
vertex}\label{faces-meeting-a-fixed-vertex}

Fix a vertex \(o\). Write \[
F_n=F_n^{\mathrm{hit}}+F_n^{\mathrm{avoid}},
\] where the two terms count faces meeting and avoiding \(o\),
respectively.

Expose all local matchings except \(V_o\). The components of the exposed
boundary graph are:

\begin{itemize}
\tightlist
\item
  cycles, which are precisely the faces avoiding \(o\);
\item
  paths whose endpoints are flags at \(o\).
\end{itemize}

The latter paths induce a perfect matching \(P\) on the \(2d\) flags at
\(o\). Conditional on this exposure, \(V_o\) is still uniform on
\(\mathcal H_d\), and the faces meeting \(o\) correspond exactly to the
components of \(P\cup V_o\).

Expose \(V_o\) pair by pair. After \(r\le d-2\) pairs have been exposed,
select any free flag \(x\). At most one possible mate closes a component
of \(P\) together with the exposed pairs: the opposite endpoint of the
current \(P\)-path containing \(x\). By (1), the conditional probability
of completing a component is at most \[
\frac{1}{2(d-r-1)}.
\] The final forced pair completes at most one component. Every
component is counted when its last pair is exposed. Therefore \[
\boxed{
\mathbb E F_n^{\mathrm{hit}}
\le
1+\frac12 H_{d-1}
=
1+\frac12 H_{n-2}.
}
\qquad\text{(2)}
\]

The remainder of the proof bounds \(F_n^{\mathrm{avoid}}\).

\begin{center}\rule{0.5\linewidth}{0.5pt}\end{center}

\subsubsection{\texorpdfstring{3. Exploring one face while avoiding
\(o\)}{3. Exploring one face while avoiding o}}\label{exploring-one-face-while-avoiding-o}

We now use a different exposure order, analyzing one starting flag at a
time and revealing no other faces.

Choose a flag \(a_0\) on an edge not incident with \(o\), and put \[
b_0=A(a_0).
\] Start with the \(A\)-edge \(a_0b_0\), treating \(a_0\) as the active
flag and \(b_0\) as a fixed terminal.

At each step:

\begin{enumerate}
\def\labelenumi{\arabic{enumi}.}
\tightlist
\item
  reveal the \(V\)-mate \(z\) of the active flag;
\item
  if \(z=b_0\), the face closes;
\item
  otherwise, make \(A(z)\) the new active flag;
\item
  stop if the new active flag is at \(o\).
\end{enumerate}

Until closure, all exposed \(V\)-pairs lie on one path with endpoints
the active flag and \(b_0\). All other components of \(A\) together with
these exposed pairs are isolated \(A\)-edges. Thus every query is to a
previously unexposed mate.

If the trace closes after \(k\) queries, its face has length \(k\).

\paragraph{The escape probability}\label{the-escape-probability}

Suppose the process is alive, its active flag is at \(u\ne o\), and
\(r_u\) pairs have been exposed at \(u\).

Neither flag of incidence \((u,uo)\) has been used. Otherwise the trace
would already have visited \(o\). Consequently those two flags form an
isolated \(J_u\)-path, different from the path containing the active
flag. Both are allowed choices in the local exposure lemma.

In particular, there are at least two unfinished local paths, so the
forced-last-pair case cannot occur while the trace avoids \(o\). The
conditional probability of moving next to \(o\) is exactly \[
p_o=\frac{2}{2(d-r_u-1)}
=\frac{1}{n-2-r_u}.
\qquad\text{(3)}
\] For any other vertex \(w\), there are at most two flags at \(u\)
whose \(A\)-partners lie at \(w\). Some might already be unavailable,
and choosing the terminal may instead close the face. Thus the
conditional probability \(p_w\) of making \(w\) the next active vertex
satisfies \[
p_w\le p_o.
\qquad\text{(4)}
\]

Let \(T\) be the number of queries until closure or the first visit to
\(o\). Set \[
q=1-\frac{1}{n-2}.
\] Equation (3) gives \[
\boxed{\Pr(T>t)\le q^t.}
\qquad\text{(5)}
\]

\paragraph{A geometric bound on repeated
visits}\label{a-geometric-bound-on-repeated-visits}

Fix \(w\ne o\). While the trace is alive, let \(N_w(t)\) count its
active visits to \(w\), including the current active position and the
initial position if applicable.

Define \[
M_t=
\begin{cases}
2^{N_w(t)},&t<T,\\
0,&t\ge T.
\end{cases}
\] Let \(p_{\mathrm{cl}}\) denote the conditional probability of closing
at the next query. While alive, \[
\mathbb E[M_{t+1}\mid\mathcal F_t]
=
M_t(1-p_o-p_{\mathrm{cl}}+p_w)
\le M_t
\] by (4). Hence \(M_t\) is a nonnegative supermartingale, with
\(M_0\le2\).

Stop it on reaching \(L\) active visits to \(w\), or on termination.
There are at most \(n(n-1)\) queries, so bounded stopping gives \[
\boxed{
\Pr\bigl(N_w(t)\ge L\text{ for some }t<T\bigr)
\le 2^{1-L}.
}
\qquad\text{(6)}
\]

This controls local congestion throughout the entire trace, not merely
up to a predetermined length.

\begin{center}\rule{0.5\linewidth}{0.5pt}\end{center}

\subsubsection{\texorpdfstring{4. Bounding all faces avoiding
\(o\)}{4. Bounding all faces avoiding o}}\label{bounding-all-faces-avoiding-o}

For sufficiently large \(n\), choose \[
R=\left\lceil4\log_2 n\right\rceil,
\qquad
D=n-2-R>0.
\qquad\text{(7)}
\] Call an avoiding face \textbf{good} if it uses at most \(R\) local
pairs at every vertex, and \textbf{bad} otherwise.

Let \(S\) be the flags on edges not incident with \(o\). Then \[
|S|=4\binom{n-1}{2}=2(n-1)(n-2).
\qquad\text{(8)}
\]

\paragraph{\texorpdfstring{4.1. Bad avoiding faces contribute
\(o(1)\)}{4.1. Bad avoiding faces contribute o(1)}}\label{bad-avoiding-faces-contribute-o1}

If the face of a starting flag avoids \(o\) and is bad, its trace
reaches \(R+1\) active visits to some vertex before terminating. By (6)
and a union bound, \[
\Pr(\text{the starting flag lies in a bad avoiding face})
\le (n-1)2^{-R}.
\qquad\text{(9)}
\]

For a flag \(a\), let \(L(a)\) be its face length. Every length-\(k\)
face has \(2k\) flags, so the number \(B\) of bad avoiding faces is
exactly \[
B=
\sum_{a\in S}
\frac{\mathbf 1_{\{a\text{ lies in a bad avoiding face}\}}}{2L(a)}.
\] Using \(L(a)\ge1\), (8), and (9), \[
\mathbb E B
\le
(n-1)^2(n-2)2^{-R}
\le n^3\,2^{-R}
\le \frac1n.
\qquad\text{(10)}
\]

No independence between traces from different starting flags is used.

\paragraph{4.2. Closing probabilities for good
faces}\label{closing-probabilities-for-good-faces}

Fix a starting flag, and let \(w_0\) be the vertex containing its
terminal \(b_0\).

For the trace to close at query \(k\), it must:

\begin{itemize}
\tightlist
\item
  remain alive after query \(k-2\);
\item
  at query \(k-1\), move to \(w_0\);
\item
  at query \(k\), choose the particular terminal flag \(b_0\).
\end{itemize}

When a revealed history has used at most \(R\) local pairs at every
vertex, the exposure rule gives at least \(2D\) possible partners at
every relevant query. Therefore:

\begin{itemize}
\tightlist
\item
  the probability of moving next to \(w_0\) is at most \(2/(2D)=1/D\);
\item
  once there, the probability of choosing \(b_0\) is at most \(1/(2D)\).
\end{itemize}

These are bounds conditional on the revealed history, \textbf{not} on
eventual goodness. To estimate a good-face event, discard continuations
whose histories cease to satisfy the load restriction. Applying the two
bounds successively, and then (5), gives, for \(k\ge2\), \[
\boxed{
\Pr(\text{the starting flag lies in a good avoiding face of length }k)
\le
\frac{q^{k-2}}{2D^2}.
}
\qquad\text{(11)}
\]

This estimate applies to every possible face length.

\paragraph{4.3. Summation over lengths}\label{summation-over-lengths}

For \(n\ge3\), there are no faces of length one or two.

Length one is impossible because \(K_n\) has no loops. A length-two face
would have to use both sides of one edge \(uv\). At each endpoint, this
would require \(V_u\) or \(V_v\) to pair the two flags of that same
incidence. Such a pair would form a proper alternating cycle with
\(J_u\) or \(J_v\), contrary to membership in \(\mathcal H_d\) when
\(d>1\).

Let \(G_k\) count good avoiding faces of length \(k\). Rooting at flags,
(8) and (11) give \begin{equation*}
\begin{aligned}
\mathbb E G_k
&=
\frac1{2k}
\sum_{a\in S}
\Pr(a\text{ lies in such a face})\\
&\le
\frac{(n-1)(n-2)}{2D^2}\,
\frac{q^{k-2}}{k}.
\end{aligned}
\end{equation*} Together with (10), \[
\boxed{
\mathbb E F_n^{\mathrm{avoid}}
\le
\frac{(n-1)(n-2)}{2(n-2-R)^2}
\sum_{k=3}^{\infty}\frac{q^{k-2}}{k}
+\frac1n.
}
\qquad\text{(12)}
\]

The series is explicit: \begin{equation*}
\begin{aligned}
\sum_{k=3}^{\infty}\frac{q^{k-2}}{k}
&=
q^{-2}\left(-\ln(1-q)-q-\frac{q^2}{2}\right)\\
&=
q^{-2}\left(\ln(n-2)-q-\frac{q^2}{2}\right)\\
&=
\ln n-\frac32+O\!\left(\frac{\ln n}{n}\right).
\end{aligned}
\qquad\text{(13)}
\end{equation*} Also, because \(R=O(\ln n)\), \[
\frac{(n-1)(n-2)}{2(n-2-R)^2}
=
\frac12+O\!\left(\frac{\ln n}{n}\right).
\qquad\text{(14)}
\] Consequently, \[
\boxed{
\mathbb E F_n^{\mathrm{avoid}}
\le
\frac12\ln n-\frac34
+O\!\left(\frac{(\ln n)^2}{n}\right).
}
\qquad\text{(15)}
\]

\begin{center}\rule{0.5\linewidth}{0.5pt}\end{center}

\subsubsection{5. Conclusion}\label{conclusion}

From (2), \[
\mathbb E F_n^{\mathrm{hit}}
\le
1+\frac12H_{n-2}
=
\frac12\ln n+1+\frac{\gamma}{2}+O(1/n).
\] Adding (15) proves \[
\mathbb E F_n
\le
\ln n+\frac14+\frac{\gamma}{2}
+O\!\left(\frac{(\ln n)^2}{n}\right).
\]

This establishes the conjecture. The argument was stated for
sufficiently large \(n\); finitely many smaller values are absorbed into
the \(O(1)\) in the conjecture.

\paragraph{If one conditions on
non-orientability}\label{if-one-conditions-on-non-orientability}

The stipulated independent-signature model also produces orientable
embeddings with a very small probability. If ``non-orientable random
embedding'' is instead intended to mean conditioning on
non-orientability, the conclusion remains unchanged.

For a connected graph, precisely \(2^{n-1}\) signatures are orientable:
these are obtained by switching vertex orientations from an
all-orientable signature. Thus, for \(K_n\), the orientability
probability is \[
p_n=2^{n-1-\binom n2}
=2^{-(n-1)(n-2)/2}.
\] Since \(F_n\ge0\), \[
\mathbb E[F_n\mid\text{non-orientable}]
\le
\frac{\mathbb E F_n}{1-p_n},
\] which preserves the displayed asymptotic upper bound.

The long-face gap is therefore closed: faces meeting the fixed vertex
are controlled by a single local matching exposure, while avoiding faces
admit both an exponential survival bound and a geometric
local-congestion bound. No additional assumption about mixing or long
facial walks is required.

\appendix

\section{Provenance and review
status}\label{provenance-and-review-status}

\begin{warning}

\textbf{UNREFEREED MODEL OUTPUT.} This document was selected solely
because GPT-6 Astra labelled its own result
\texttt{would\_publish:\ true} and returned \texttt{proved} or
\texttt{disproved}. It has not passed the adversarial LLM referee used
for the earlier Sol campaign, has not been checked by a human
mathematician, and has not been checked for novelty. Treat every
mathematical and bibliographic claim as unverified.

\end{warning}

{\def\LTcaptype{none} % do not increment counter
\begin{longtable}[]{@{}
  >{\raggedright\arraybackslash}p{(\linewidth - 2\tabcolsep) * \real{0.5000}}
  >{\raggedright\arraybackslash}p{(\linewidth - 2\tabcolsep) * \real{0.5000}}@{}}
\toprule\noalign{}
\endhead
\bottomrule\noalign{}
\endlastfoot
Catalog id & \texttt{2211.01032\_\_03} \\
Catalog entry &
\href{https://graph-theory-ai.github.io/graph-conjectures/arxiv/2211.01032__03/}{Conjecture
9.5} \\
Source corpus & arxiv \\
Campaign leg & selected second attempts (\texttt{attacks\_retry}) \\
Source paper / entry & Random Embeddings of Graphs: The Expected Number
of Faces in Most Graph\ldots{} \\
Model verdict & \textbf{proved} (confidence: high) \\
Selection criterion & \texttt{would\_publish:\ true} and verdict
\texttt{proved} or \texttt{disproved} \\
Model & \texttt{gpt-6-astra}, reasoning effort \texttt{max},
\texttt{mode=pro}, flex service tier \\
Original artifact &
\texttt{attacks\_retry/2211.01032\_\_03/output.md} \\
Independent review & \textbf{None} \\
arXiv & \href{https://arxiv.org/abs/2211.01032}{arXiv:2211.01032} \\
Previous attempt & \texttt{attacks/2211.01032\_\_03}: gpt-5.6-sol
verdict \texttt{partial} \\
Model's caveats & The proof uses exactly the stated uniform-rotation,
independent-fair-signature model; literature priority has not been
checked. \\
\end{longtable}
}

The generated Markdown and LaTeX sources for this PDF are stored under
\texttt{to\_review\_astra/src/2211.01032\_\_03/}.

\end{document}
